\section{Overview of the Processing Pipeline} \label{sec:pipeline-overview}

After dataset cleaning, the analysis pipeline began with a comprehensive \textbf{preprocessing} stage.
This included \textit{timestamp localization} to ensure all \gls{eeg} recordings were temporally aligned, accounting for time zone differences and \gls{dst} changes.
\textit{Artifacts} in the \gls{eeg} signals, such as extreme amplitude events or technical noise, were systematically detected and removed to prevent contamination of downstream analyses.
The continuous \gls{eeg} data were then \textit{segmented} into short, fixed-length windows (\textit{segments}) and longer, non-overlapping windows (\textit{clips}), with careful labeling of intervals relative to seizure events (e.g., preictal, interictal).
The dataset was \textit{partitioned} into training and test sets, ensuring sufficient representation of both periictal and interictal periods for robust model development and evaluation.

\textbf{Feature extraction} was performed on each segment, yielding a set of quantitative descriptors that captured both temporal and spectral properties of the \gls{eeg}.
These included signal variance, \gls{acfw}, Pearson correlation between channels, and absolute band powers across standard frequency bands.
Some features were used as direct inputs to machine learning models to align with literature, while all were retained for the complementary cycle analysis.
The extraction process was designed to maximize the interpretability and physiological relevance of the features, supporting both predictive modeling and the investigation of cyclical patterns in seizure risk.

Two main \textbf{model architectures} were employed: a \gls{cnn} that operated on raw \gls{eeg} segments, and an ensemble of \glspl{fb-mlp} that utilized the extracted features.
The \gls{cnn} architecture was adapted to the specific characteristics of the dataset, including its sampling rate and channel configuration, and was regularized to prevent overfitting.
The \gls{fb-mlp} ensemble leveraged z-score normalized features and aggregated predictions across multiple independently initialized models.
Both models were trained using only preictal and interictal segments, with careful handling of class imbalance and standardized training procedures to ensure comparability.

\textbf{Model evaluation} was conducted using both standard machine learning metrics and \glspl{ebm} tailored to the clinical context of seizure prediction.
The latter included measures like \gls{tifw} and \gls{eb} sensitivity, which reflected the real-world utility of the models in predicting seizures while minimizing unnecessary alarms.
Thresholds for issuing warnings were optimized to balance \gls{eb} sensitivity and \gls{eb} specificity, and all evaluations were performed separately on training and test sets to assess generalization.

To quantify the \textbf{co-occurrence of false predictions}, the \gls{cope} was computed per patient between the predictions of both models.

To further understand the temporal dynamics of seizure risk, \textbf{cycle extraction} was performed on the feature time series.
This involved gap filling, outlier handling, and bandpass filtering to isolate multidien (multi-day) oscillatory components.
The relationship between cyclical features and event timing was quantified using circular statistics, including \gls{plv} and the Rayleigh test for phase preference.
A non-parametric permutation test was used to compare the phase distributions of seizures and model-predicted \glspl{fp}, with multiple testing correction applied to control the \gls{fdr}.

Finally, the \textbf{propensity evaluation} framework integrated these analyses to determine whether a model’s \glspl{fp} were meaningfully aligned with periods of increased seizure risk.
Qualification criteria were defined for each patient, feature, and model, requiring significant phase locking and predictive performance.
Only when these criteria were met were the phase relationships between \glspl{fp} and seizures interpreted as evidence of the model capturing cyclical seizure propensity, supporting the overall objective of understanding and predicting seizure risk in a clinically relevant manner.
